\section{Misceláneos}
% - % - % - % - % - % - % - % - % - % - % - % - % - % - % - % -
% - % - % - % - % - % - % - % - % - % - % - % - % - % - % - % -
\subsection{XOR Basis}


\algoinfo{d}{d}{} 
\begin{minted}{cpp}
int basis[d];// basis[i]: v\inBase cuyo lsb es i
int sz;       // size actual de la base (= rango r)
void insertVector(int mask) {
    for( int i = 0; i < d; i++ ){
        if( !is_on(mask,i) )continue; // continue if lsb(mask)!=i
        if( !basis[i] ){ // no hay vector en la base con ese bit
            basis[i] = mask;
            ++sz;
            return;
        }
        mask ^= basis[i]; // si ya hay, quitarle ese bit a mask
    }
}
\end{minted}
\begin{extended}[] \vspace{-2mm}
\begin{description}
\item[Teorema de Rango--Nulidad.] 
Sea $\varphi:W\to V$ una transformación lineal con $W$ de dimensión finita. Entonces
$
\dim W=\dim\ker \varphi+\dim\operatorname{im}\varphi.
$

\item[Primer Teorema de Isomorfismo]
Utilizando $W / \ker \varphi \cong \operatorname{im}\varphi$, los tamaños de los conjuntos satisfacen:
$
|W| = |\ker \varphi| \cdot |\operatorname{im} \varphi|.
$
\item[Cardinalidad de un Espacio Vectorial Finito]
Si $V$ es un espacio vectorial de dimensión finita sobre un cuerpo finito $K$, la cardinalidad de $V$ cumple que:
$
|V| = |K|^{n} \text{, si } V \cong K^n.
$

\item[Teorema de las Fibras]
Sea $\varphi: W \to V$ una transformación lineal entre espacios vectoriales. Para todo vector $v \in \operatorname{im}\varphi$, la fibra (preimagen) $\varphi^{-1}(v)$ es una clase lateral del núcleo, dada por:
\[
\varphi^{-1}(v) = w_0 + \ker \varphi = \{ w_0 + k \mid k \in \ker \varphi \}
\]
donde $w_0 \in W$ es cualquier solución particular tal que $\varphi(w_0) = v$. 

Como consecuencia, existe una biyección entre $\ker \varphi$ y cualquier fibra no vacía $\varphi^{-1}(v)$, por lo que la cardinalidad de cada fibra coincide exactamente con la del núcleo:
\[
|\varphi^{-1}(v)| = |\ker \varphi| \quad \forall \, v \in \operatorname{im}\varphi.
\]
\end{description}
\end{extended}


\textbf{Resultado}
Sean $v_1,\dots,v_n\in\mathbb F_2^d$ y sea $V=\operatorname{span}(v_1,\dots,v_n)$
con $r=\dim V$. Se quiere contar, para $x\in\mathbb F_2^d$, cuántos
$S\subseteq[n]$ cumplen $\bigoplus_{i\in S}v_i=x$. (Si $x\notin V$, la respuesta es 0.)
\[
\#\Bigl\{S\subseteq[n] : \bigoplus_{i\in S}v_i=x\Bigr\}=
\begin{cases}
2^{\,n-r} & \text{si } x\in V,\\
0 & \text{si } x\notin V.
\end{cases}
\]



